\subsection{从数据库 ID 到列表顺序的物理位置}

列表排序使分配给同一 IVF 质心的向量连续存储。输入行的标识为 $i\in\{0,\ldots,N-1\}$；粗分配使用计算得到的质心分数，其数学目标为
\begin{equation}
 a(i)=\operatorname*{arg\,min}_{0\leq\ell<n_{\rm list}}
       \|\Pi x_i-\mu_\ell\|_2^2.
 \label{eq:4-assignment}
\end{equation}
实现采用上文说明的数值分配路径。主码和残差编码不减去 $\mu_{a(i)}$；列表归属独立于这些重建，限制后续搜索空间。

构建器将分配结果与输入 ID 一同排序。令 $\rho(p)$ 为排序后物理位置 $p$ 上的输入 ID，$o_\ell$ 为列表 $\ell$ 的起点。完成的数组满足
\begin{equation}
 0=o_0\leq o_1\leq\cdots\leq o_{n_{\rm list}}=N,
 \qquad p\in[o_\ell,o_{\ell+1})\Longrightarrow a(\rho(p))=\ell.
 \label{eq:4-offsets}
\end{equation}
lower-bound 搜索找到各列表的首个位置，因此相同偏移表示空列表：
\[
 o_\ell=o_{\ell+1}\Longleftrightarrow\text{list }\ell\text{ is empty}.
\]
位置在完成的索引中保持固定，重建后可能改变。所有数组使用同一个 $\rho(p)$，ID 通过 $p\mapsto\rho(p)$ 映射回输入行。两遍 add 暂时构造 $\rho^{-1}$，将新编码直接写入目标位置，称为 scatter。主码并列使用物理位置，最终并列使用数据库 ID。

\par\medskip\noindent\begin{minipage}{\linewidth}
\captionsetup{hypcap=false}
\centering
% Source-level schematic: v69 gather/repack and separate allocations.
% All addresses below are byte offsets relative to each allocation's base.
\begin{tikzpicture}[x=1cm,y=1cm,font=\footnotesize\linespread{1}\selectfont,
 cell/.style={draw=ThesisNavy!50,minimum width=.88cm,minimum height=.48cm,inner sep=1pt},
 panel/.style={draw=ThesisNavy!40,rounded corners=2pt,fill=ThesisNavy!2},
 link/.style={-{Stealth[length=4pt]},ThesisNavy!80,semithick}]
\node[anchor=west,font=\small\bfseries] at (0,1.0) {List ranges and physical positions (schematic, $N=8$)};
\node[anchor=west] at (0,.13) {$o=[0,3,3,8]$};
\node[anchor=west] at (0,-.32) {list 1: $[3,3)=\varnothing$};
\foreach \p in {0,...,7} {
 \node[cell,fill=ThesisNavy!5] (pos\p) at (4.4+.88*\p,-.3) {$\p$};
}
\draw[decorate,decoration={brace,amplitude=3pt}] (3.96,.12)--(6.60,.12)
 node[midway,above=4pt] {list 0: $[0,3)$};
\draw[decorate,decoration={brace,amplitude=3pt}] (6.60,.12)--(11.00,.12)
 node[midway,above=4pt] {list 2: $[3,8)$};
\node[anchor=west,align=left] at (11.55,-.05) {Each range indexes\\all four buffers.};
\node[draw=ThesisNavy,fill=ThesisNavy!12,rounded corners=2pt,
 minimum height=.5cm] (chosen) at (1.25,-1.02) {example $p=3$};
\draw[link] (pos3.south)--(7.04,-.72)-|(chosen.north);
\node[anchor=west] at (4.0,-1.02) {Separate allocations; links below denote an index join.};

% ID allocation.
\draw[panel] (1.05,-1.53) rectangle (14.85,-2.57);
\node[anchor=west,align=left] at (1.3,-2.04) {\textbf{IDs}\\$\mathrm{IDs}[p]=\rho(p)$};
\foreach \p/\id in {0/5,1/1,2/7,3/0,4/6,5/3,6/4,7/2} {
 \node[cell,fill=white] at (4.4+.88*\p,-2.04) {$\id$};
}
\node[cell,draw=ThesisNavy,fill=ThesisNavy!17] (id3) at (7.04,-2.04) {$0$};
\node[anchor=west,align=left] at (11.5,-2.04) {int32: $4p$\\$p=3:\ 12$ bytes};

% Primary allocation, word-group-major. The example shows M >= 8.
\draw[panel] (1.05,-2.88) rectangle (14.85,-5.45);
\node[anchor=west,align=left] at (1.3,-3.62) {\textbf{Primary words}\\group-major\\(shown: $M\geq8$)};
\foreach \j/\yy in {0/-3.36,1/-3.96} {
 \foreach \p in {0,...,7} {
  \node[cell,fill=white] at (4.4+.88*\p,\yy) {$\mathsf w_{\j,\p}$};
 }
 \node[cell,draw=ThesisNavy,fill=ThesisNavy!17] (word\j) at (7.04,\yy) {$\mathsf w_{\j,3}$};
}
\node[anchor=west,align=left] at (11.5,-3.65) {$4(jN+p)$\\$p=3,N=8:$\\$j=0:\ 12$ bytes\\$j=1:\ 44$ bytes};
\draw[link] (4.0,-4.43)--(10.95,-4.43)
 node[midway,below=2pt] {fixed $j$: adjacent $p$ give consecutive 32-bit words};
\node[anchor=west] at (1.3,-5.12)
 {One word: $Z[4j,p]$, $Z[4j+1,p]$, $Z[4j+2,p]$, $Z[4j+3,p]$ (one candidate).};

% Residual allocation; each displayed cell represents a whole row.
\draw[panel] (1.05,-5.76) rectangle (14.85,-7.53);
\node[anchor=west,align=left] at (1.3,-6.33) {\textbf{Residual codes}\\row-major (8-bit)};
\foreach \p in {0,...,7} {
 \node[cell,fill=white] at (4.4+.88*\p,-6.33) {row $\p$};
}
\node[cell,draw=ThesisNavy,fill=ThesisNavy!17] (res3) at (7.04,-6.33) {row 3};
\node[anchor=west,align=left] at (11.5,-6.33) {$pb_r+t$\\$p=3:\ 3b_r+t$};
\node[anchor=west] at (1.3,-7.10) {row 3 contains $b_r$ consecutive bytes:};
\node[draw=ThesisNavy!50,fill=white,minimum height=.45cm,inner sep=3pt]
 (expanded) at (9.05,-7.10) {$R[3,0]\quad R[3,1]\quad\cdots\quad R[3,b_r-1]$};
\draw[link] (res3.south)--(expanded.north west);

% Correction allocation.
\draw[panel] (1.05,-7.84) rectangle (14.85,-9.12);
\node[anchor=west,align=left] at (1.3,-8.48) {\textbf{Correction}\\$\mathrm{correction}[p]$\\$=\gamma[p]$};
\foreach \p in {0,...,7} {
 \node[cell,fill=white] at (4.4+.88*\p,-8.36) {$\gamma[\p]$};
}
\node[cell,draw=ThesisNavy,fill=ThesisNavy!17] (corr3) at (7.04,-8.36) {$\gamma[3]$};
\node[anchor=west,align=left] at (11.5,-8.36) {FP32: $4p$\\$p=3:\ 12$ bytes};

% One physical-position key points into four independent buffers.
\draw[ThesisNavy!80,semithick] (chosen.west)--(.35,-1.02)--(.35,-7.69);
\draw[link] (.35,-1.38)-|(id3.north);
\draw[link] (.35,-2.73)-|(word0.north);
\draw[link] (.35,-5.61)-|(res3.north);
\draw[link] (.35,-7.69)-|(corr3.north);
\node[anchor=west,align=left] at (0,-9.65)
 {All offsets are relative to each buffer's base. Refinement fetches survivor rows only;\\
 survivor positions can be scattered. Highlighted entries share $p=3$, not an interleaved allocation.};
\end{tikzpicture}

\captionof{figure}[GPU 索引的物理布局]{GPU 索引的物理布局.}
\label{fig:4-layout}
\end{minipage}\par\medskip

图中示例取 $N=8$，offsets $o=[0,3,3,8]$ 使 list~1 为空。浅色列标出 $p=3$：四个独立数组均以相同物理位置索引，右侧 byte 偏移相对各自数组的起点。固定 word group $j$ 时，相邻候选对应连续的 32 位 word，使扫描相邻候选的线程可以合并访存；八位残差配置下，每行存放同一向量的 $b_r$ 个编码 bytes。

\subsection{精确地址与表示宽度}

主码表示按子空间组织候选。对主码 byte $Z[m,p]$，地址为
\[
 \operatorname{offset}_{\rm byte}(Z[m,p])=mN+p.
\]
最终打包将同一向量的四个子空间 byte 合为一个 word：
\begin{equation}
 \mathsf w[j,p]=\sum_{t=0}^{3}2^{8t}Z[4j+t,p],\qquad
 \operatorname{addr}_{\rm byte}(\mathsf w[j,p])=4(jN+p).
 \label{eq:4-primary-address}
\end{equation}
这里 $j=0,\ldots,M/4-1$。$\text{word-layout build}\Longrightarrow M\bmod4=0$；其余 $M$ 被拒绝，独立的 byte-layout 编译则保留 byte 表示。每个数组占 $NM$ bytes，重打包暂时同时持有两份副本，成本为 $2NM$。一遍构建将已有编码重排（gather）到 vector-major 的 $[N,M]$ 数组，转置并重打包，逐次释放旧副本；两遍构建直接 scatter 到 subspace-major 存储，再进行相同重打包。搜索时的 lane 映射见 \Cref{sec:5-3}。

残差码保持 row-major，每向量占 $b_r=\lceil dB_r/8\rceil$ bytes。八位配置中，行 $p$ 的坐标 $t$ 位于 byte 偏移 $pb_r+t$；四位路径中，它位于 byte $pb_r+\lfloor t/2\rfloor$，偶数坐标在低四位（low nibble），奇数坐标在高四位（high nibble）。每个物理行有四个 byte 的校正项，存储
\begin{equation}
 \gamma[ p ]=2\langle\widehat p(x_{\rho(p)}),
                              \widehat r(x_{\rho(p)})\rangle.
 \label{eq:4-correction}
\end{equation}
gather 或 scatter 时，校正项必须与该行编码一同移动。按列表扫描随后只为保留候选读取残差。精排使用压缩重建；原始 FP32 重排序还需增加 $4Nd$ 个设备 bytes。第 7 章评估其召回率。
